\item \textbf{Requerimientos no funcionales}

\begin{itemize}
        \item Asociados a requerimientos funcionales

Una máquina física pertenece exclusivamente a una sola sucursal.

El costo inicial de la tarjeta de acceso está fijado obligatoriamente en \$20.00 MXN, requiriendo además una recarga inicial mínima de \$10.00 MXN.

Las recargas de saldo y los costos de todas las partidas deben ser estrictamente múltiplos de 10.

Para obtener la categoría VIP, el jugador debe registrar 10 visitas en un mismo mes.

Una "visita" se contabiliza al usar la tarjeta en una sucursal durante un día determinado, sin importar la cantidad de partidas jugadas ese día.

La categoría VIP otorga un 10\% adicional sobre los puntos obtenidos y requiere 10 visitas mensuales para mantenerse activa.

El canje de premios requiere un mínimo de 20 puntos.

Los premios cuestan en puntos el doble de su valor monetario aproximado.

Los premios están estrictamente divididos por rangos de edad: Infantil (3 a 12 años), Juvenil (13 a 17 años) y Adulto (18 años en adelante).  
Las categorías del inventario de premios se dividen en Bajos (20 a 1,000 puntos), Medios (1,001 a 3,999 puntos) y Grandes (4,000 a más de 10,000 puntos).


        \item Asociados a requerimientos no funcionales

Para el prototipo, la persistencia de datos debe realizarse obligatoriamente guardando y leyendo la información en archivos .csv.

El prototipo debe operar mediante un menú de interacción.

El código del prototipo debe incluir manejo de excepciones para notificar al usuario si ocurre un error durante la ejecución.

Los campos de entrada numéricos deben ser validados para aceptar única y exclusivamente números.
\end{itemize}
